% ================================================
% 文件: 01_术语表.tex
% 章节: 附录A 术语表
% 所属部分: 附录
% 版本: v1.0
% ================================================

\chapter{术语表}
\label{chap:glossary}

\section{本章目录}
\begin{itemize}
    \item 第1节 无线电基础术语
    \item 第2节 接收机术语
    \item 第3节 发射机术语
    \item 第4节 系统与设备术语
\end{itemize}

\section{无线电基础术语}

\begin{description}
    \item[RF] 射频 (Radio Frequency)
    \item[IF] 中频 (Intermediate Frequency)
    \item[AF] 音频 (Audio Frequency)
    \item[Hz] 赫兹，频率的国际单位制（SI）单位，1 Hz = 1 周/秒
    \item[kHz] 千赫兹，$10^3$ Hz
    \item[MHz] 兆赫兹，$10^6$ Hz
    \item[周 / 周率] 频率的旧式中文称呼（20 世纪中期），"千周"即 kHz
    \item[MW] 中波 (Medium Wave)，广播频段 526.5–1606.5 kHz（ITU 1 区）
    \item[AM] 调幅 (Amplitude Modulation)
    \item[FM] 调频 (Frequency Modulation)
    \item[调制] 将低频信号加载到高频载波上的过程
    \item[解调] 从已调载波中提取原始信号的过程
    \item[谐振] 电路对特定频率产生最大响应的现象
    \item[调谐] 调整电路使其与特定频率谐振的过程
    \item[ANT] 天线 (Antenna)，用于发射或接收电磁波
    \item[GND] 信号/工作接地 (Ground)
    \item[PE] 保护接地 (Protective Earthing)
\end{description}

\section{接收机术语}

\begin{description}
    \item[灵敏度] 接收机检测微弱信号的能力
    \item[选择性] 接收机区分不同频率信号的能力
    \item[超外差] 将射频信号转换为固定中频进行处理的架构
    \item[检波] 从调幅信号中提取音频信号的过程
    \item[AGC] 自动增益控制 (Automatic Gain Control)
    \item[S/N] 信噪比 (Signal-to-Noise Ratio)
\end{description}

\section{发射机术语}

\begin{description}
    \item[功率放大器] 放大射频信号功率的电路
    \item[振荡器] 产生高频信号的电路
    \item[调制器] 将音频信号加载到载波上的电路
    \item[PA] 功率放大器 (Power Amplifier)
    \item[效率] 输出功率与输入功率的比值
\end{description}

\section{系统与设备术语}

\begin{description}
    \item[LED] 发光二极管 (Light Emitting Diode)
    \item[RFID] 射频识别 (Radio Frequency Identification)
    \item[天线] 发射或接收无线电波的装置
    \item[传输线] 传输高频信号的导线
    \item[连接器] 连接不同设备或线缆的器件
\end{description}